% part: intuitionistic-logic
% chapter: propositions-as-types
% section: normalization

\documentclass[../../../include/open-logic-section]{subfiles}

\begin{document}

\olfileid{int}{pty}{nor}

\olsection{عادي کول}

په دې برخه کښې ښيو چې د راکمولو د يوه ترتيب په وسيله هر
صحيح ټايپ‌لرونکے ثبوت عادي ثبوت ته راکمېږي؛ دې ته
\emph{د عادي کولو خاصيت} وايي۔ د بيانونو-په-ټايپونو سمون
کاروو: ښيو چې هر صحيح ټايپ‌لرونکے ثبوت ترم عادي بڼې ته
راکمېږي؛ بيا د فطري استنتاج د !!{derivation}s عادي کول ترې
راځي۔ \textbf{د ژباړې د نسخې سپيناوی (OLCMP-201):} دا
د صحيح ټايپ‌لرونکو ترمونو دعوه ده، د خامو ثبوت ترمونو د
هرې نحوي څرګندونې نۀ ده۔ د دې تجربوي باب درې محاسبوي
راکمول د هغه په خپل حد کښې دي؛ د ځاے بدلولو او ساده کولو
د پاتې بدلونونو بشپړه کتنه له دې دعوې نۀ ثابتېږي۔

لومړے د ترمونو د پېچلتيا د اندازې دپاره ځينې تابعې تعريفوو۔
د !!a{formula}s \emph{اوږدوالے}~$\len{!A}$ داسې تعريفېږي:
\begin{align*}
  \len{p} &= 0\\
  \len{\lfalse} &= 0\\
  \len{!A \land !B} &= \len{!A} + \len{!B} + 1\\
  \len{!A \lor !B} &= \len{!A} + \len{!B} + 1 \\
  \len{!A \lif !B} &= \len{!A} + \len{!B} + 1.
\end{align*}
د رېډکس~$M$ پېچلتيا د هغه په \emph{پرېکون رتبه}
~$\cutrank{M}$ اندازه کوو:
\begin{align*}
  \cutrank{(\lambd[\typeof{x}{!A}][\typeof{N}{!B}])Q} &=
  \len{!A} + \len{!B} + 1 \\
  \cutrank{\ande{i}{\andi{\typeof{M}{!A}}{\typeof{N}{!B}}}} &=
  \len{!A} + \len{!B} + 1 \\
  \cutrank{\ore{\ori{i}{!A_{i}}{\typeof{M}{!A_i}}}
    {\typeof{x_1}{!A_1}}{\typeof{N_1}{!C}}
    {\typeof{x_2}{!A_2}}{\typeof{N_2}{!C}}} &=
  \len{!A_1} + \len{!A_2} + 1
\end{align*}
د ثبوت ترم پېچلتيا په هغه کښې د تر ټولو پېچلي رېډکس رتبه ده؛
که عادي وي، $0$ ده:
\[
  \maxrank{M} = \max (\{0\} \cup \{\cutrank{N} | N \text{ د } M \text{
    فرعي ترم او رېډکس دے}\})
\]
\textbf{د ژباړې د نسخې سمون (OLCMP-202):} د تش ټايپ
اوږدوالے هم صفر ټاکل شو؛ سرچينه ئې نۀ ټاکي۔ د يا د
پرېکون په وروستۍ رتبه کښې د هماغو دوو ښودل شوو ټايپونو
شاخصونه راغلل، د بې‌تعريفه دوو نومونو پر ځاے۔ د صفر
ورګډول د بې‌رېډکسه ترم د نثري صفر قاعده په رسمي اعظمي
اندازه کښې هم څرګندوي۔ دا رتبه د همدې تعريف د اوږدوالي
اندازه ده، د مخکنيو بابونو د فورمول د ژوروالي عين اندازه نۀ ده۔

\begin{lem}\ollabel{lem:subst}
  که $\Subst{M}{\typeof{N}{!A}}{\typeof{x}{!A}}$ رېډکس وي او
  $M \not\ident x$ وي، نو له لاندې حالاتو څخه يو پېښېږي:
  \begin{enumerate}
  \item خپله $M$ رېډکس دے، يا
  \item $M$ د $\ande{i}{x}$ په بڼه او $N$ د
    $\andi{P_1}{P_2}$ په بڼه دے۔
  \item $M$ د $\ore{x}{x_1}{P_1}{x_2}{P_2}$ په بڼه او $N$
    د $\ori{i}{}{Q}$ په بڼه دے۔
  \item $M$ د $x Q$ په بڼه او $N$ د
    $\lambd[x][P]$ په بڼه دے۔
  \end{enumerate}
  په لومړي حالت کښې $\cutrank{\Subst{M}{N}{x}} = \cutrank{M}$؛
  په نورو حالاتو کښې $\cutrank{\Subst{M}{N}{x}} = \len{!A}$۔
\end{lem}
\begin{proof}
  پر~$M$ د استقرا ثبوت دے۔
  \begin{enumerate}
  \item که $M$ يوازې متغير $y$ وي او $y \not\ident x$ وي، نو
    $\Subst{y}{N}{x}$ هم $y$ دے، نو رېډکس نۀ دے۔
  \item که $M$ د $\andi{N_1}{N_2}$، $\lambd[x][N]$ يا
    $\ori{i}{!A}{N}$ په بڼه وي، نو
    $\Subst{M}{\typeof{N}{!A}}{\typeof{x}{!A}}$ هم په هماغه
    بڼه دے، نو رېډکس نۀ دے۔
  \item که $M$ د $\ande{i}{P}$ په بڼه وي، دوه حالتونه ګورو۔
    \begin{enumerate}
      \item که $P$ د $\andi{P_1}{P_2}$ په بڼه وي، نو
        $M \ident \ande{i}{\andi{P_1}{P_2}}$ رېډکس دے؛ په
        څرګند ډول
        \[
        \Subst{M}{N}{x} \ident
        \ande{i}{\andi{\Subst{P_1}{N}{x}}{\Subst{P_2}{N}{x}}}
        \]
        هم رېډکس دے۔ د پرېکون رتبې برابرې دي۔
      \item که $P$ يوازې متغير وي، بايد $x$ وي، څو تعويض
        رېډکس شي؛ او $N$ بايد د $\andi{P_1}{P_2}$ په بڼه وي۔
        اوس
        \[
        \Subst{M}{N}{x} \ident \Subst{\ande{i}{x}}{\andi{P_1}{P_2}}{x},
        \]
        وګورئ؛ دا $\ande{i}{\andi{P_1}{P_2}}$ دے۔ د پرېکون
        رتبه ئې د تعويض کېدونکي جوړې د ټايپ اوږدوالے، يعنې
        $\len{!A}$، دے۔
    \end{enumerate}
  \end{enumerate}
  د $\ore{N}{x_1}{N_1}{x_2}{N_2}$ او $P Q$ حالتونه ورته دي۔
  \textbf{د ژباړې د نسخې سمون (OLCMP-203):} د يا ايستلو
  په بيان کښې د تعويض کېدونکي متغير نوم په لومړي ځای کښې
  راغے؛ عددي شاخص هلته غلط و۔ د اوږدوالي له معادلې يو
  بې‌جوړه قوس لرې شو۔ متغير په خپله رېډکس نۀ دے، نو د
  هغه د پرېکون بې‌تعريفه رتبه د جوړې د ټايپ اوږدوالي ته
  اصلاح شوه؛ نور رسمي ترمونه عين دي۔
\end{proof}

\begin{lem}
  که $M$ ترم $M'$ ته انقباض وکړي، او د $M$ د هر مناسب
  رېډکس فرعي ترم $N$ دپاره $\cutrank{M} > \cutrank{N}$ وي،
  نو $\cutrank{M} > \maxrank{M'}$۔
\end{lem}

\begin{proof}
  د حالاتو له مخې ثبوت دے۔
  \begin{enumerate}
  \item که $M$ د $\ande{i}{\andi{M_1}{M_2}}$ په بڼه وي، نو
    $M'$، $M_i$ دے؛ د $M_i$ هر فرعي ترم د $M$ مناسب
    فرعي ترم هم دے، نو دعوه صحيح ده۔
  \item که $M$ د
    $(\lambd[\typeof{x}{!A}][N])\typeof{Q}{!A}$ په بڼه وي، نو
    $M'$، $\Subst{N}{\typeof{Q}{!A}}{\typeof{x}{!A}}$ دے۔ په
    $M'$ کښې يو رېډکس وګورئ۔ يا په $N$ کښې ورته رېډکس
    د برابرې پرېکون رتبې سره شته، چې د فرض له مخې له
    $\cutrank{M}$ څخه کمه ده؛ يا رتبه $\len{!A}$ ده، چې
    د تعريف له مخې له $\cutrank{(\lambd[x^{!A}][N])Q}$ څخه
    کمه ده۔ د آرګومېنټ په نقل شوو برخو کښې د پخوا شته
    رېډکسونو حالت هم شته: هغوی د اصلي ټول رېډکس مناسب
    فرعي رېډکسونه وو، نو د فرض له مخې ئې رتبې ټيټې دي۔
  \item که $M$ د
    \[
    \ore{\ori{i}{}{\typeof{N}{!A_i}}}
        {\typeof{x_1}{!A_1}}{\typeof{N_1}{!C}}
        {\typeof{x_2}{!A_2}}{\typeof{N_2}{!C}},
    \]
    په بڼه وي، نو $M' \ident \Subst{N_i}{N}{\typeof{x_i}{!A_i}}$۔
    په~$M'$ کښې يو رېډکس وګورئ۔ يا په $N_i$ کښې ورته
    رېډکس د برابرې پرېکون رتبې سره شته، چې د فرض له مخې
    له $\cutrank{M}$ څخه کمه ده؛ يا رتبه $\len{!A_i}$ ده،
    چې د تعريف له مخې له
    $\cutrank{\ore{\ori{i}{}{\typeof{N}{!A_i}}}
      {\typeof{x_1}{!A_1}}{\typeof{N_1}{!C}}
      {\typeof{x_2}{!A_2}}{\typeof{N_2}{!C}}}$ څخه کمه ده۔
    د ورننوتونکي آرګومېنټ په نقلونو کښې پخواني رېډکسونه هم
    د اصلي رېډکس مناسب فرعي ترمونه دي، نو د فرض له مخې
    ټيټې رتبې لري۔
  \end{enumerate}
  \textbf{د ژباړې د نسخې سمون (OLCMP-204):} سرچينې د
  بدنې پخواني او د تعويض په پوله نوي جوړ شوي رېډکسونه
  ياد کړي، خو د آرګومېنټ له نقلونو سره راغلي پخواني
  رېډکسونه ئې پرېښي وو؛ دغه اړين حالت اوس څرګند دے۔
\end{proof}

\begin{thm}
  ټول صحيح ټايپ‌لرونکي ثبوت ترمونه عادي بڼې ته راکمېږي؛
  د دې باب په اړوند محاسبوي معنا ټول !!{derivation}s عادي
  !!{derivation}s ته راکمېږي۔
\end{thm}

\begin{proof}
  دويمه پايله له لومړۍ راځي۔ د لومړۍ دپاره د
  $m=\maxrank{M}$ پر اندازه بشپړه استقرا کوو، چې $M$
  صحيح ټايپ‌لرونکے ثبوت ترم دے۔
  \begin{enumerate}

  \item که $m=0$ وي، $M$ له مخکې عادي دے۔
  \item
    که نه وي، پر~$n$ استقرا کوو؛ دا په $M$ کښې د هغو
    رېډکسونو شمېر دے چې د پرېکون رتبه ئې~$m$ ده۔
    \begin{enumerate}
    \item که $n=1$ وي، داسې هر رېډکس $N$ غوره کړئ چې
      $m = \cutrank{N} > \cutrank{P}$ د هر مناسب فرعي ترم
      $P$ دپاره چې رېډکس هم وي، صحيح وي۔ داسې رېډکس شته،
      ځکه هر ترم يوازې متناهي شمېر فرعي ترمونه لري۔

      د $N$ راکمې شوې پايلې ته $N'$ وايو۔ د لمې له مخې
      $\maxrank{N'} < \maxrank{N}$؛ نو سرچينه استدلال کوي چې
      $n$، يعنې د $\cutrank=m$ لرونکو رېډکسونو شمېر، کمېږي۔
      له دې سره $m$ هم (په $1$ يا زيات) کمېږي، او د~$m$
      استقرايي فرضيه کارولے شو۔
    \item د استقرا په ګام کښې $n>1$ فرض کړئ۔ بهير ورته
      دے، خو $n$ يوازې يوه مثبته شمېره ته کمېږي، نو $m$
      نۀ بدلېږي۔ د~$n$ استقرايي فرضيه کاروو۔
    \end{enumerate}
\end{enumerate}
\textbf{د ژباړې د نسخې د ثبوت حد (OLCMP-205):} سرچينه
د ټاکلي فرعي رېډکس د پايلې ټيټه رتبه ښيي؛ د ټول ترم د
شمېر د کمېدو دپاره بايد په چاپېريال کښې د ممکنو نوو
دباندنيو رېډکسونو رتبې هم وکتل شي۔ دغه عمومي ګام دلته
نۀ دے بشپړ شوے؛ يو بديل بشپړ ثبوت نۀ دے اختراع شوے۔
دا ټاکلې تګلاره د عادي بڼې د وجود دعوه ده؛ د هر ترتيب
درېدنه له همدې استقرا نۀ ثابتېږي۔
\end{proof}

د ټايپ‌لرونکو $\lambd$-ترمونو عادي کېدل وايي چې هر داسې
ترم \emph{عادي بڼه لري}۔ که د $\lambd$-ترمونو عادي بڼې
ته $\beta$-راکمول د محاسبې اجرا او عادي بڼه ئې د محاسبې
ارزښت وګڼو، نو په ټايپ‌لرونکي $\lambd$-حساب کښې د ټاکلي مناسب ترتيب په وسيله محاسبه درېږي
او ارزښت ورکوي۔ د ټايپ ساتنه هم د ارزښت صحيح ټايپ
تضمينوي۔ مثلاً که $N$ ټايپ $!A \lif !B$ ولري (يعنې د
ټايپ~$!A$ آرګومېنټونو دپاره تابع ده او د ټايپ~$!B$
ارزښتونه ورکوي)، او په د ټايپ~$!A$ ترم~$M$ (داخلې)
ئې تطبيق کړو، نو $(NM)$ ترم $N'$ (وت) ته راکمېږي او
دغه وت په يقيني ډول ټايپ~$!B$ لري۔

سرچينه يو بل قوي خاصيت هم بيانوي: په ټايپ‌لرونکي
$\lambd$-حساب کښې ترمونه يوازې پر فرعي ترمونو د
$\beta$-راکمول د يوه ترتيب له مخې عادي بڼې ته نۀ رسېږي؛
بلکې د $\beta$-راکمول \emph{هر} ترتيب بالاخره په عادي
بڼه کښې درېږي۔ دغه خاصيت \emph{قوي عادي کول} بلل کېږي۔
دا وايي چې يوازې د محاسبې د درېدنې او ارزښت د ورکولو
يوه تضمين شوې لاره نشته، بلکې د محاسبې \emph{هره}
اجرا بالاخره ارزښت ورکوي۔ د OLCMP-205 د حد له مخې د
دې اضافي قوي دعوې ثبوت دلته نۀ دے وړاندې شوے؛ د هغې
رښتينولي يوازې له مخکني کمزوري دليل څخه نۀ اخلو۔

څنګه پوهېږو چې د محاسبې د اجرا بېلابېل ترتيبونه
\emph{يو شان} ارزښت ورکوي؟ تر اوسه د ټايپ ساتنه يوازې
وايي چې ټول ورکړي ارزښتونه يو ټايپ لري۔ ټايپ‌لرونکے
$\lambd$-حساب يو بل مهم خاصيت لري: دا \emph{هم‌پايلے}
يا \emph{چرچ--روسر} دے۔ که $N \red N_1$ او $N \red N_2$
وي، نو داسې ترم~$N'$ شته چې $N_1 \red N'$ او $N_2 \red N'$
وي۔ په ياد ولرئ چې $N \red N'$ معنا لري چې $N'$ له پيلني ترم
څخه د صفر يا زياتو $\redone$-راکمول ګامونو په وسيله راځي۔
که $N_1$ عادي بڼه وي، نو يوازې هغه ترم $N'$ چې
$N_1 \red N'$ ورته صحيح وي، خپله $N_1$ دے (د صفر
$\redone$-راکمول ګامونو په وسيله)۔ نو که $N_1$ او $N_2$
دواړه عادي بڼې وي، $N_1 = N' = N_2$۔ په بل عبارت،
د $\red$ د قوي عادي کېدو له خاصيته هر ترتيب بالاخره
عادي بڼې ته رسېږي؛ او د $\red$ د چرچ--روسر له خاصيته
هرې دوې عادي بڼې برابرې دي۔ نو په ټايپ‌لرونکي
$\lambd$-حساب کښې محاسبه تل درېږي او ارزښت (عادي بڼه)
ئې د اجرا له ترتيب څخه خپلواک دے۔ دلته دغه پايله د
يادو دواړو خاصيتونو پر رښتينوالي ولاړه ده؛ د مخکني
نيمګړي دليل په وسيله د دواړو بشپړ ثبوت نۀ ادعا کېږي۔

د يوې اړيکې د هم‌پايلتوب ثبوت عموماً ګران وي۔ خو
لاندې کمزورے شرط په اسانه ښودلے شو:

\begin{prop}[د چرچ--روسر کمزورے خاصيت]
که $N \redone N_1$ او $N \redone N_2$ وي، نو داسې ترم
$N'$ شته چې $N_1 \red N'$ او $N_2 \red N'$ وي۔
\end{prop}

\begin{proof}
  دوه ترمونه $N_1$ او $N_2$ په ترتيب سره د $N$ د رېډکس
  $M_1$ يا $M_2$ په راکمې شوې پايله بدلولو څخه راځي۔
  رېډکسونه $M_1$ او $M_2$ د $N$ فرعي ترمونه دي، نو د
  ترم په ونه کښې جزوي متقاطع نۀ دي: يا يو د بل فرعي
  ترم دے، يا د $N$ په بېلو نۀ پوښل کېدونکو ځایونو کښې
  واقع دي۔ وروستے حالت وګورئ: $N$ د $N(M_1,M_2)$ په
  بڼه دے۔ د $M_1$ پر ځاے د هغه راکمې پايلې $M_1'$
  اېښودل $N_1$ ورکوي، چې $N(M_1',M_2)$ دے؛ او د
  $M_2$ پر ځاے د هغه راکمې پايلې $M_2'$ اېښودل $N_2$
  ورکوي، چې $N(M_1,M_2')$ دے۔ دواړه $N(M_1',M_2')$
  ته راکمېږي۔ د بل حالت دپاره فرض کړئ چې $M_2$ د~$M_1$
  فرعي ترم دے (بل حالت متناظر دے)، او د $M_1$ د رېډکس
  د ډول له مخې حالتونه بېل کړئ۔ مثلاً که
  $M_1 \ident \proj{1}{\pair{O_1}{O_2}}$ وي، نو $O_1$
  ته راکمېږي۔ که $M_2$ د $O_1$ فرعي ترم وي، د $M_2$
  بدلول $O_1'$ ورکوي۔ نو لومړے د $M_1$ او ورپسې د
  $M_2$ بدلول $O_1'$ ورکوي۔ خو لومړے د $M_2$ بدلول
  $\proj{1}{\pair{O_1'}{O_2}}$ ورکوي، او دا هم~$O_1'$
  ته راکمېږي۔ \textbf{د ژباړې د نسخې سمون (OLCMP-206):}
  د سرچينې وروستے دويم جز غلط و؛ د لومړي پروجکشن پايله
  تازه لومړے جز دے۔ سرچينه د دننيو رېډکسونو نور ډولونه
  يوازې يادوي؛ د نورو ډولونو بشپړ حالتونه دلته نۀ دي
  وړاندې شوي، نو دغه بېلګه د ټولو حالتونو بشپړ ثبوت
  نۀ ګڼل کېږي۔
\end{proof}

\begin{lem}[د نيومن لمه]
که $\red$ قوي عادي کېدونکے وي او د چرچ--روسر کمزورے
خاصيت ولري، نو هم‌پايلے دے۔
\end{lem}

\begin{proof}
  څرنګه چې $\red$ قوي عادي کېدونکے دے، هره تر پايه
  وړل شوې د راکمولو لړۍ عادي بڼې ته رسېږي۔ د دې
  فرض لاندې عادي بڼې يکتا دي هله او يوازې هله چې
  $\red$ هم‌پايلے وي۔ فرض کړئ چې يو ترم $N$ دوې
  بېلې عادي بڼې $N'$ او $N''$ ولري، چې د لاندې
  راکمولو لړيو پايلې دي:
  \begin{align*}
  N & = N_1' \redone N_2' \redone \dots \redone N_k' = N' \text{ او}\\
  N & = N_1'' \redone N_2'' \redone \dots \redone N_l'' = N''
  \end{align*}
  (د راکمولو لړۍ بايد اوږدوالے $>0$ ولري، ځکه که نه
  وي، $N$ به له مخکې عادي بڼه وه۔)

  که $N_2' = N_2''$ وي، همدغه سمدستي پايله دوې
  بېلې عادي بڼې ($N'$ او $N''$) لري۔ که
  $N_2' \neq N_2''$ وي، د چرچ--روسر د کمزوري خاصيت
  له مخې داسې $N'''$ شته چې $N_2' \red N'''$ او
  $N_2'' \red N'''$ وي۔ دغه ګډه پايله د درېدنې په
  وسيله يوې عادي بڼې $K$ ته ورسوئ؛ نو
  $N_2' \red K$ او $N_2'' \red K$۔ څرنګه چې
  $N' \neq N''$ دي، يا $K \neq N'$ يا $K \neq N''$۔
  نو يا $N_2'$ يا $N_2''$ دوې بېلې عادي بڼې لري۔
  ومو ښودله چې که $N$ دوې بېلې عادي بڼې ولري، داسې
  $N_1$ شته چې $N \redone N_1$ وي او $N_1$ دوې
  بېلې عادي بڼې ولري۔ دغه ګام تکرارول
  $N \redone N_1 \redone N_2 \redone \dots$ ورکوي،
  خو دا ناممکن دے، ځکه $\redone$ قوي عادي کېدونکے دے۔
  \textbf{د ژباړې د نسخې سمون (OLCMP-207):} سرچينې د
  لومړي راکم شوي ترم پر ځاے د لړ هماغه پيل بيا اخيستے
  و۔ اوس دويم لړي غړي پرتله کېږي، څو ګام په رښتيا
  سمدستي پايلې ته لاړ شي۔ ګډه پايله پخپله عادي نۀ
  فرض کېږي؛ لومړے ئې عادي بڼه اخلو، څو د دوو بېلو
  عادي بڼو استدلال صحيح وي۔
\end{proof}

\end{document}
